% =====================================================================
% GUÍA DE ESTILO Y REVISIÓN
% Este capítulo solo se imprime en modo borrador (main.tex lo incluye
% dentro de \ifdraft). Desaparece automáticamente al pasar a \draftfalse,
% por lo que NO debes borrarlo ni rellenarlo: es una guía para ti.
% No aparece en el índice (no lleva \addcontentsline).
% =====================================================================
\chapter*{Guía de estilo y revisión (no aparece en la versión final)}
\markboth{Guía de estilo y revisión (no aparece en la versión final)}{}

Este capítulo es una guía para ti: solo se imprime en modo borrador y
desaparece de forma automática al compilar con \texttt{\textbackslash
draftfalse}. Léelo antes de empezar a redactar y vuelve a él durante la
revisión final. Una memoria con buen contenido pero mal escrita pierde
puntos fáciles; las convenciones siguientes evitan los errores más
frecuentes.

\section*{Registro y voz}

Elige una voz y mantenla en toda la memoria: la forma impersonal o pasiva
refleja (\enquote{se ha desarrollado}, \enquote{se propone}) o la primera
persona del plural (\enquote{hemos desarrollado}). Ambas son aceptables;
mezclarlas no. La única excepción son los agradecimientos, que van en
primera persona del singular.

Usa el tiempo pasado (pretérito perfecto) para el trabajo realizado
(\enquote{se ha implementado un módulo de autenticación}) y el presente
para describir lo que el sistema es o hace (\enquote{el sistema valida
cada petición}). Reserva el futuro y el condicional para el trabajo
futuro. Evita los coloquialismos y muletillas (\enquote{un montón},
\enquote{básicamente}, \enquote{a día de hoy}) y no te dirijas al lector
en segunda persona.

\section*{Vocabulario técnico}

Un concepto, un término: si llamas \enquote{usuario} a quien usa la
aplicación, no alternes con \enquote{actor} o \enquote{cliente} para
referirte a lo mismo. Haz una pasada final de búsqueda para comprobar la
consistencia. Los anglicismos sin traducción asentada van en cursiva con
\texttt{\textbackslash textit}: \textit{framework}, \textit{endpoint},
\textit{backend}. Si existe un término español habitual (prueba,
despliegue, biblioteca, servidor), úsalo, y úsalo siempre igual.

Los comandos, el código, las rutas y los identificadores van en
\texttt{\textbackslash texttt}: \texttt{main.tex}, \texttt{POST
/api/login}, \texttt{docker compose up}. Todo acrónimo se desarrolla en
su primera mención en el texto ---\enquote{Ingeniería del Software Basada
en Evidencias (EBSE)}--- y a partir de entonces se usa solo la sigla;
además, debe quedar recogido en la lista de acrónimos.

\section*{Prosa antes que viñetas}

Las listas-guía de esta plantilla son andamiaje: indican qué contenido
debe aparecer, no cómo debe presentarse. El texto final debe ser prosa
argumentada, con párrafos conectados mediante conectores (\enquote{Sin
embargo}, \enquote{Por tanto}, \enquote{En consecuencia}). Una sección no
puede ser solo una sucesión de viñetas; reserva las listas para
enumeraciones reales (requisitos, objetivos, pasos de un procedimiento).
Las citas textuales van siempre con \texttt{\textbackslash
enquote\{\dots\}} (el paquete \texttt{csquotes} ya está cargado), nunca
con comillas tecleadas a mano.

\section*{Figuras y tablas}

Toda figura y toda tabla se referencia desde el texto \textbf{antes} de
aparecer (\enquote{como muestra la Tabla~3.1\dots}, usando
\texttt{Tabla\textasciitilde\textbackslash ref\{\dots\}}) y se interpreta:
no basta con insertarla, hay que decir al lector qué debe ver en ella.
Los pies deben ser autoexplicativos, de modo que se entiendan sin leer el
cuerpo del texto (\enquote{Comparación del tiempo medio de respuesta (ms)
de las tres alternativas evaluadas} en lugar de \enquote{Resultados}).

\section*{Checklist de revisión final}

\begin{itemize}[itemsep=0.3em]
  \item[$\square$] Corrector ortográfico y gramatical pasado a todo el
    documento (el integrado en Overleaf y, como complemento, LanguageTool
    o similar).
  \item[$\square$] Lectura completa del PDF de principio a fin, mejor
    impresa o en voz alta: delata las frases que no fluyen.
  \item[$\square$] Terminología consistente: el mismo término para el
    mismo concepto en toda la memoria; anglicismos en cursiva; código y
    comandos en \texttt{\textbackslash texttt}.
  \item[$\square$] Acrónimos desarrollados en su primera mención y
    recogidos en la lista de acrónimos (sin entradas sobrantes).
  \item[$\square$] Compilado con \texttt{\textbackslash draftfalse} y
    revisado el log de compilación: ningún aviso \enquote{Placeholder sin
    rellenar} pendiente.
  \item[$\square$] Ningún título de sección ni caption de tabla/figura
    sigue siendo un placeholder (en modo final dejarían entradas vacías
    en el índice o pies vacíos).
  \item[$\square$] Toda entrada de la bibliografía se cita en el texto y
    toda afirmación tomada de una fuente externa lleva su cita.
  \item[$\square$] Todas las figuras y tablas se referencian desde el
    texto y tienen pie autoexplicativo.
  \item[$\square$] Metadatos del PDF completos: bloque \enquote{DATOS DEL
    TRABAJO} de \texttt{main.tex} relleno (título, autor/a, titulación,
    tutores, ciudad y fecha).
\end{itemize}
